% !TeX root = ../guide
\chapter{Notepad++: A Lightweight Text Editor for Advanced Workflows}\label{ch:notepadplusplus}
\chaptermark{Notepad++}
	\section{Introduction}\label{sec:NPPIntroduction}
		While dedicated IDEs like TeXstudio are highly recommended for writing the bulk of your thesis, there is undeniable value in having a fast, lightweight text editor in your toolkit. 
		Notepad++ is a free, open-source source code editor that supports several languages, including \LaTeX.
		
		Unlike TeXstudio, Notepad++ will not hold your hand with integrated PDF viewers or built-in mathematical symbol libraries. 
		However, its speed, robust search capabilities, and flexibility make it an excellent secondary tool for managing raw code, editing configuration files, or cleaning up messy bibliography data.
	
	\section{Key Features of Notepad++}
		Notepad++ strips away the heavy graphical interfaces of full IDEs, focusing entirely on efficient text manipulation.

		\subsection{Speed and Low Resource Usage}
			Because it is a plain text editor at its core, Notepad++ opens almost instantly, even when loading massive text files or raw data sets. 
			This makes it ideal for quick edits when you do not want to wait for a full IDE to initialize.

		\subsection{Advanced Search and Replace}
			Perhaps the most powerful feature of Notepad++ is its \enquote{Find in Files} capability. 
			If you realize you have used the wrong terminology or need to change a variable name across all twenty chapters of your thesis, Notepad++ can search your entire project folder and replace the text in one click. 
			It also fully supports Regular Expressions (Regex) for highly complex search patterns.
		
		\subsection{Syntax Highlighting}
			Out of the box, Notepad++ recognizes \path{.tex} and \path{.bib} files and will apply basic syntax highlighting to your \LaTeX\ commands, making the raw code much easier to read than in standard Windows Notepad.
		
		\subsection{Plugin Ecosystem}
			Through the built-in Plugin Admin, you can extend Notepad++'s functionality. 
			Plugins like \textit{NppExec} allow you to send commands directly to your computer's command line, meaning you can set up a custom shortcut to compile your \LaTeX\ document directly from the editor.
	
	\section{Getting Started with Notepad++}
		Setting up Notepad++ is quick and straightforward. 
		
		\subsection{Installation}
			Notepad++ is available exclusively for Windows. 
			To install the software, visit the official website: \url{https://notepad-plus-plus.org/}.
			Navigate to the `Download' section, select the latest release, and download the installer. 
			Follow the standard Windows installation prompts.
		
		\subsection{The Main Interface}
			Upon launching, you will see a clean, tabbed interface as shown in \Cref{fig:NotepadMain}.
			\begin{figure}[htbp]
				\centering
				\includegraphics[width=0.8\linewidth]{example-image}
				\caption{Notepad++ Main Interface showing a \LaTeX\ document with syntax highlighting.}
				\label{fig:NotepadMain}
			\end{figure}
		
		\subsection{Compiling \LaTeX\ from Notepad++ (Advanced)}
			If you wish to use Notepad++ as your primary editor, you will need a way to compile your document into a PDF. This requires the \textit{NppExec} plugin:
			\begin{enumerate}
				\item Go to `Plugins' $\rightarrow$ `Plugins Admin'.
				\item Search for `NppExec', check the box, and click `Install'. Notepad++ will restart.
				\item Press `F6' to open the Execute dialog.
				\item Enter the following commands to compile your document:\par
					\texttt{cd "\$(CURRENT\_DIRECTORY)"})\par
					\texttt{latexmk -pdf -silent -synctex=1 "\$(FILE\_NAME)"})
				\item Save this script for future use.
			\end{enumerate}
			\note{%
				For most thesis writers, it is far easier to use TeXstudio for compiling and reserve Notepad++ for quick text edits.}
	
	\section{Notepad++ in a Thesis Workflow}
		You might wonder why you need Notepad++ if you already have TeXstudio. 
		Here are a few scenarios where a lightweight text editor shines:
		
		\subsection{Editing Style (.sty) and Class (.cls) Files}
			If you need to make a minor tweak to the uAlberta thesis template's underlying code (such as adjusting a margin or fixing a title page layout), opening the \texttt{.cls} file in Notepad++ is often cleaner and safer than loading it into your main IDE.
		
		\subsection{Cleaning up Bib\TeX\ Files}
			Sometimes, exporting references from online databases results in messy \texttt{.bib} files with broken characters or formatting errors. 
			Opening your bibliography database in Notepad++ allows you to view the raw text without any software trying to \enquote{interpret} it, making it easier to spot and delete corrupted lines.
		
		\subsection{Managing External Scripts}
			If your research involves writing Python or other programming scripts for data analysis or generating complex TikZ plots, Notepad++ allows you to edit those external scripts side-by-side with your \LaTeX\ code in a lightweight environment.